\documentclass[12pt,a4paper]{amsart}

\usepackage{amsmath,amssymb,amsthm}
\usepackage{mathtools}
\usepackage[ngerman]{babel}
\usepackage{hyperref}
\usepackage{geometry}
\geometry{margin=2.5cm}

% -------------------------------------------------------
% Theorem-Umgebungen
% -------------------------------------------------------
\newtheorem{theorem}{Satz}[section]
\newtheorem{lemma}[theorem]{Lemma}
\newtheorem{corollary}[theorem]{Korollar}
\newtheorem{proposition}[theorem]{Proposition}
\theoremstyle{definition}
\newtheorem{definition}[theorem]{Definition}
\newtheorem{remark}[theorem]{Bemerkung}
\newtheorem{example}[theorem]{Beispiel}

% -------------------------------------------------------
% Eigene Befehle
% -------------------------------------------------------
\newcommand{\N}{\mathbb{N}}
\newcommand{\Z}{\mathbb{Z}}
\newcommand{\Q}{\mathbb{Q}}
\newcommand{\R}{\mathbb{R}}
\newcommand{\C}{\mathbb{C}}
\newcommand{\Fp}{\mathbb{F}_p}
\DeclareMathOperator{\ord}{ord}

% -------------------------------------------------------
\begin{document}
% -------------------------------------------------------

\title{Elliptische Kurven \"uber \texorpdfstring{$\Q$}{Q}:\\
       Weierstra\ss-Form, Gruppengesetz, Torsion und der Satz von Mordell--Weil}

\author{Michael Fuhrmann}
\address{Specialist Mathematics Project, Build 84}
\email{specialist-maths@localhost}
\date{\today}

\begin{abstract}
Eine \emph{elliptische Kurve} \"uber $\Q$ ist eine glatte projektive Kurve
vom Geschlecht~$1$ mit einem ausgew\"ahlten rationalen Punkt, typischerweise
in \emph{kurzer Weierstra\ss-Form} $y^2 = x^3 + ax + b$ mit $a, b \in \Q$
und nicht-verschwindender Diskriminante $\Delta = -16(4a^3 + 27b^2) \ne 0$.
Wir entwickeln die Theorie systematisch:
die algebraische Gruppenstruktur durch das \emph{Sekanten-Tangenten-Gesetz},
die Diskriminante und die \emph{$j$-Invariante},
die \emph{Torsionsuntergruppe} $E(\Q)_{\mathrm{tors}}$ (klassifiziert durch
Mazurs Satz in $15$ m\"ogliche Typen)
und den \emph{Satz von Mordell--Weil}
$E(\Q) \cong E(\Q)_{\mathrm{tors}} \oplus \Z^r$,
der besagt, dass die Gruppe der rationalen Punkte endlich erzeugt ist.
\end{abstract}

\maketitle
\tableofcontents

% ================================================================
\section{Einleitung}
% ================================================================

Elliptische Kurven sind zentral f\"ur Zahlentheorie, Kryptographie und
algebraische Geometrie.
Sie sind die einfachsten nicht-trivialen Kurven aus arithmetischer Sicht:
Rationale Punkte auf Kurven vom Geschlecht~$0$ sind leicht parametrisiert
(Kegelschnitte), w\"ahrend Geschlecht~$\ge 2$ durch Faltings' Satz
(endlich viele rationale Punkte) beherrscht wird.
Geschlecht~$1$ -- der elliptische Fall -- ist reich:
Die rationalen Punkte bilden eine abelsche Gruppe,
und diese Gruppe zu bestimmen ist ein tiefes arithmetisches Problem.

% ================================================================
\section{Weierstra\ss-Form und grundlegende Definitionen}
% ================================================================

\begin{definition}[Kurze Weierstra\ss-Form]
\label{def:weierstrass}
Eine \emph{elliptische Kurve} \"uber $\Q$ in \emph{kurzer Weierstra\ss-Form} ist
eine ebene kubische Kurve
\[
  E: \quad y^2 = x^3 + ax + b
\]
mit $a, b \in \Q$ und \emph{Diskriminante}
\[
  \Delta \;=\; -16\,(4a^3 + 27b^2) \;\ne\; 0.
\]
Die Bedingung $\Delta \ne 0$ sichert die Glattheit von $E$.
Zusammen mit dem \emph{Punkt im Unendlichen} $\mathcal{O}$ ist $E$ eine
glatte projektive Kurve vom Geschlecht~$1$.
\end{definition}

\begin{definition}[$j$-Invariante]
\label{def:j_invariant}
Die \emph{$j$-Invariante} von $E: y^2 = x^3 + ax + b$ ist
\[
  j(E) \;=\; 1728\,\frac{4a^3}{4a^3 + 27b^2}.
\]
\end{definition}

\begin{example}[Wichtige Sonderkurven]
\label{ex:special_curves}
\begin{itemize}
  \item $j = 1728$: Die Kurve $E\colon y^2 = x^3 - x$ hat $a = -1$, $b = 0$
    und $j(E) = 1728 \cdot \frac{4(-1)^3}{4(-1)^3 + 27\cdot 0} = 1728$.
    Sie ist die kongruente Zahlkurve f\"ur $n=1$ und hat komplexe Multiplikation
    durch $\Z[i]$.
  \item $j = 0$: Die Kurve $E\colon y^2 = x^3 + 1$ hat $a = 0$, $b = 1$
    und $j(E) = 1728 \cdot \frac{0}{0 + 27} = 0$.
    Sie hat komplexe Multiplikation durch $\Z[\omega]$ mit $\omega = e^{2\pi i/3}$.
\end{itemize}
Die $j$-Werte $0$ und $1728$ sind die einzigen Kurven mit zus\"atzlichen
Automorphismen (Ordnung $6$ bzw.\ $4$).
\end{example}

\begin{theorem}[Klassifikation durch $j$]
\label{thm:j_classification}
Zwei elliptische Kurven $E$ und $E'$ \"uber $\overline{\Q}$ sind genau dann
isomorph als elliptische Kurven, wenn $j(E) = j(E')$.
\"Uber $\Q$ bestimmt die $j$-Invariante die Kurve bis auf einen
\emph{Twist} (Verdrehung).
\end{theorem}

\begin{proof}[Beweisskizze]
Ein Isomorphismus $E \to E'$ von Weierstra\ss-Kurven wird durch
$(x,y) \mapsto (u^2 x + r, u^3 y + u^2 sx + t)$ gegeben.
Unter dieser Abbildung transformieren sich $a \mapsto u^4 a'$ und
$b \mapsto u^6 b'$, woraus $j = j'$ folgt.
Umgekehrt gibt es zu jedem $j$-Wert ein Weierstra\ss-Modell; vgl.\ \cite{Silverman1986}, Kap.~III.
\end{proof}

% ================================================================
\section{Das Gruppengesetz: Sekanten-Tangenten-Regel}
% ================================================================

\begin{theorem}[Gruppengesetz auf $E$]
\label{thm:group_law}
Die Menge $E(\Q) = \{(x,y) \in \Q^2 : y^2 = x^3+ax+b\} \cup \{\mathcal{O}\}$
bildet eine abelsche Gruppe mit neutralem Element $\mathcal{O}$ und
folgender Additionsregel:

\begin{enumerate}
  \item[\textup{(i)}] $\mathcal{O}$ ist das neutrale Element.
  \item[\textup{(ii)}] Inverses: $-(x,y) = (x,-y)$.
  \item[\textup{(iii)}] Addition ($P \ne \pm Q$): Die Gerade durch $P = (x_1,y_1)$
    und $Q = (x_2,y_2)$ (bzw.\ die Tangente an $E$ im Punkt $P$ falls $P=Q$)
    trifft $E$ in einem dritten Punkt $R'$, und $P+Q = -R'$.
\end{enumerate}

Explizit: F\"ur $P = (x_1,y_1) \ne \pm Q = (x_2,y_2)$:
\begin{align}
  \lambda &=
  \begin{cases}
    \dfrac{y_2 - y_1}{x_2 - x_1} & \text{falls } x_1 \ne x_2
      \quad\text{(Sekante)},\\[6pt]
    \dfrac{3x_1^2 + a}{2y_1}     & \text{falls } P = Q
      \quad\text{(Tangente)},
  \end{cases}
  \label{eq:lambda} \\[4pt]
  x_3 &= \lambda^2 - x_1 - x_2, \label{eq:x3} \\
  y_3 &= \lambda(x_1 - x_3) - y_1. \label{eq:y3}
\end{align}
Dann gilt $P + Q = (x_3, y_3)$.
\end{theorem}

\begin{proof}
Wir verifizieren den Sekanten-Fall ($x_1 \ne x_2$).
Die Gerade $y = \lambda(x-x_1)+y_1$ trifft $y^2 = x^3+ax+b$ in den Punkten
mit $x$-Koordinaten, die das kubische Polynom
\[
  (\lambda(x-x_1)+y_1)^2 - x^3 - ax - b = 0
\]
erf\"ullen.
Nach Vi\`etes Formeln ist die Summe der drei Wurzeln gleich $\lambda^2$
(aus dem $x^2$-Koeffizienten).
Da $x_1$ und $x_2$ zwei Wurzeln sind, gilt $x_3 = \lambda^2 - x_1 - x_2$.
Einsetzen in die Geradengleichung liefert $y_3$.

Die Assoziativit\"at kann algebraisch (direkte Rechnung) oder geometrisch
(B\'ezout-Satz f\"ur Kubiken durch $9$ Punkte) bewiesen werden;
vgl.\ \cite{Silverman1986}, Kapitel~III.
\end{proof}

% ================================================================
\section{Die Torsionsuntergruppe}
% ================================================================

\begin{definition}[Torsionspunkte]
\label{def:torsion}
Ein Punkt $P \in E(\Q)$ hei\ss t \emph{Torsionspunkt}, falls $nP = \mathcal{O}$
f\"ur ein $n \ge 1$ gilt. Die \emph{Torsionsuntergruppe} ist
\[
  E(\Q)_{\mathrm{tors}} \;=\; \{P \in E(\Q) : nP = \mathcal{O}
  \text{ f\"ur ein } n \ge 1\}.
\]
\end{definition}

\begin{theorem}[Mazurs Torsionssatz, 1977]
\label{thm:mazur}
Sei $E$ eine elliptische Kurve \"uber $\Q$. Dann ist $E(\Q)_{\mathrm{tors}}$
eine der folgenden $15$ Gruppen:
\begin{align*}
  &\Z/n\Z \quad (n = 1, 2, 3, 4, 5, 6, 7, 8, 9, 10, 12), \\
  &\Z/2\Z \times \Z/2n\Z \quad (n = 1, 2, 3, 4).
\end{align*}
\end{theorem}

\begin{proof}[Beweisskizze]
Mazurs urspr\"unglicher Beweis \cite{Mazur1977} verwendet die Theorie der
Modulkurven $X_1(N)$ und Eisenstein-Ideale.
Schritte: (1) F\"ur eine Primzahl $\ell$ wird die $\ell$-Torsion durch
$X_1(\ell)$ kontrolliert, die f\"ur $\ell \ge 11$ Geschlecht $\ge 2$ hat
(Riemann--Hurwitz). (2) Nach Faltings' Satz (Geschlecht $\ge 2$) hat
$X_1(\ell)(\Q)$ nur endlich viele Punkte. (3) Explizite Rechnung schliesst
die restlichen F\"alle aus.
\end{proof}

\begin{example}[Torsionsbeispiele]
\begin{itemize}
  \item $E: y^2 = x^3 - x$ hat Torsion $\Z/2\Z \times \Z/2\Z$
    (drei $2$-Torsionspunkte: $(0,0)$, $(1,0)$, $(-1,0)$).
  \item $E: y^2 = x^3 - x^2$ ist \emph{singulär} (Knoten im Ursprung,
    $\Delta = 0$); sie ist \emph{keine} elliptische Kurve.
  \item $E: y^2 = x^3 + 1$ hat Torsion $\Z/6\Z$.
  \item $E: y^2 + y = x^3 - x^2 - 10x - 10$ hat Torsion $\Z/12\Z$
    (maximale zyklische Ordnung).
\end{itemize}
\end{example}

\begin{theorem}[Satz von Nagell--Lutz]
\label{thm:nagell_lutz}
Sei $E: y^2 = x^3 + ax + b$ mit $a, b \in \Z$.
Ist $(x_0, y_0) \in E(\Q)$ ein Torsionspunkt, so gilt $x_0, y_0 \in \Z$,
und entweder $y_0 = 0$ (Punkt der Ordnung~$2$) oder $y_0^2 \mid 4a^3 + 27b^2$.
\end{theorem}

\begin{proof}[Beweisskizze]
Folgt aus der Theorie formaler Gruppen und Reduktion modulo Primzahlen.
Ein elementarer Beweis steht in \cite{Silverman1986}, Kapitel~VIII.
\end{proof}

% ================================================================
\section{Der Satz von Mordell--Weil}
% ================================================================

\begin{theorem}[Satz von Mordell--Weil]
\label{thm:mordell_weil}
F\"ur jede elliptische Kurve $E$ \"uber $\Q$ ist $E(\Q)$ eine
endlich erzeugte abelsche Gruppe.
Nach dem Struktursatz f\"ur endlich erzeugte abelsche Gruppen gilt:
\begin{equation}
  \label{eq:mw}
  E(\Q) \;\cong\; E(\Q)_{\mathrm{tors}} \;\oplus\; \Z^r,
\end{equation}
wobei $r \ge 0$ der \emph{Rang} von $E$ \"uber $\Q$ ist.
\end{theorem}

\begin{proof}[Beweisskizze]
Mordell (1922) bewies dies f\"ur $\Q$; Weil (1928) verallgemeinerte es auf
Zahlk\"orper.
Der Beweis hat zwei Hauptteile:

\textbf{Teil 1: Schwacher Mordell--Weil-Satz.}
Es gibt $n \ge 2$ mit $E(\Q)/nE(\Q)$ endlich.
(F\"ur $n=2$: Die $2$-Abstieg-Abbildung ist injektiv in eine endliche
Gruppe, die mit dem Selmer-Raum zusammenh\"angt.)

\textbf{Teil 2: Abstieg.}
Mithilfe einer H\"ohenfunktion $h: E(\Q) \to \R_{\ge 0}$ mit
\begin{align}
  h(P+Q) + h(P-Q) &= 2h(P) + 2h(Q) + O(1), \\
  h(nP) &= n^2 h(P) + O(1),
\end{align}
zeigt man, dass jede Nebenklasse von $nE(\Q)$ in $E(\Q)$ einen
Repr\"asentanten beschr\"ankter H\"ohe enth\"alt und Mengen beschr\"ankter
H\"ohe endlich sind.

Details in \cite{Silverman1986}, Kapitel~VIII.
\end{proof}

\begin{remark}[Das Rang-Problem]
Der \emph{Rang} $r$ in~\eqref{eq:mw} ist das geheimnisvollste Invariante
von $E(\Q)$. Es gibt keinen bekannten Algorithmus, der garantiert $r$
berechnet.
Die Birch--Swinnerton-Dyer-Vermutung (Paper~27) sagt vorher:
$r = \ord_{s=1} L(E,s)$.

Bekannte Rekorde: Kurven mit Rang $\ge 28$ sind bekannt (Elkies 2006).
Ob die R\"ange unbeschr\"ankt sind, ist unbekannt.
\end{remark}

% ================================================================
\section{H\"ohen und die kanonische H\"ohe}
% ================================================================

\begin{definition}[Na\"ive H\"ohe]
\label{def:naive_height}
F\"ur $P = (x,y) \in E(\Q)$ mit $x = p/q$ (vollst\"andig gek\"urzt) ist
die \emph{na\"ive H\"ohe} $H(P) = \max(|p|,|q|)$ und die
\emph{logarithmische na\"ive H\"ohe} $h(P) = \log H(P)$.
\end{definition}

\begin{theorem}[Kanonische (N\'eron--Tate-)H\"ohe]
\label{thm:neron_tate}
Es gibt eine eindeutige Funktion $\hat h: E(\Q) \to \R_{\ge 0}$,
die \emph{kanonische H\"ohe}, mit:
\begin{enumerate}
  \item[\textup{(i)}] $\hat h(P) = h(P) + O(1)$.
  \item[\textup{(ii)}] $\hat h(nP) = n^2 \hat h(P)$ (exakt quadratisch).
  \item[\textup{(iii)}] $\hat h(P) \ge 0$, mit Gleichheit genau f\"ur
    Torsionspunkte.
\end{enumerate}
Die \emph{N\'eron--Tate-Bilinearform}
$\langle P, Q \rangle = \hat h(P+Q) - \hat h(P) - \hat h(Q)$
macht $E(\Q)/E(\Q)_{\mathrm{tors}} \otimes \R$ zu einem euklidischen Raum.
\end{theorem}

% ================================================================
\section{Reduktion modulo Primzahlen}
% ================================================================

\begin{definition}[Gute und schlechte Reduktion]
\label{def:reduction}
F\"ur eine Primzahl $p$ ist die \emph{Reduktion} $\tilde E$ modulo $p$
die Kurve $y^2 = x^3 + \bar a x + \bar b$ \"uber $\Fp$.
\begin{itemize}
  \item \emph{Gute Reduktion} bei $p$: $p \nmid \Delta$, d.\,h.\
    $\tilde E$ ist glatt.
  \item \emph{Multiplikative Reduktion} bei $p$: $\tilde E$ hat einen Knoten.
  \item \emph{Additive Reduktion} bei $p$: $\tilde E$ hat eine Spitze.
\end{itemize}
Der \emph{F\"uhrer} $N_E = \prod_p p^{f_p}$ fasst das schlechte
Reduktionsverhalten zusammen.
\end{definition}

\begin{theorem}[N\'eron--Ogg--Shafarevich]\label{thm:nos}
  Sei $\ell$ eine Primzahl mit $\ell \ne p$.  Dann hat $E$ gute
  Reduktion bei $p$ genau dann, wenn die $\ell$-adische Darstellung
  $\rho_\ell\colon \mathrm{Gal}(\overline{\Q}/\Q) \to \mathrm{GL}_2(\Z_\ell)$
  bei $p$ unverzweigt ist (d.\,h.\ die Tr\"agheitsgruppe $I_p$ trivial
  auf $T_\ell(E)$ wirkt).
\end{theorem}

% ================================================================
\section{Schlussfolgerung}
% ================================================================

Wir haben die Grundlagen elliptischer Kurven \"uber $\Q$ entwickelt:
die glatte Kubik in Weierstra\ss-Form, die $j$-Invariante,
die abelsche Gruppenstruktur durch die Sekanten-Tangenten-Regel,
die vollst\"andige Klassifikation der Torsion durch Mazurs Satz (15 Typen)
und den Satz von Mordell--Weil mit Struktur
$E(\Q)_{\mathrm{tors}} \oplus \Z^r$.

Der Rang $r$ -- die Anzahl unabh\"angiger Erzeuger unendlicher Ordnung --
ist das zentrale R\"atsel:
Er wird durch die $L$-Funktion $L(E,s)$ (Paper~26) und die
Birch--Swinnerton-Dyer-Vermutung (Paper~27) vorhergesagt,
bleibt aber im Allgemeinen unberechenbar.

% ================================================================
\begin{thebibliography}{10}

\bibitem{Silverman1986}
J.~H.~Silverman,
\emph{The Arithmetic of Elliptic Curves},
Graduate Texts in Mathematics~106,
Springer, New York, 1986.

\bibitem{Cassels1991}
J.~W.~S.~Cassels,
\emph{Lectures on Elliptic Curves},
London Mathematical Society Student Texts~24,
Cambridge University Press, 1991.

\bibitem{Mazur1977}
B.~Mazur,
Modular curves and the Eisenstein ideal,
\emph{Inst.\ Hautes \'Etudes Sci.\ Publ.\ Math.}\ \textbf{47} (1977), 33--186.

\bibitem{Washington2008}
L.~C.~Washington,
\emph{Elliptic Curves: Number Theory and Cryptography}, 2.~Aufl.,
CRC Press, 2008.

\bibitem{Cornell1997}
G.~Cornell, J.~H.~Silverman und G.~Stevens (Hrsg.),
\emph{Modular Forms and Fermat's Last Theorem},
Springer, New York, 1997.

\end{thebibliography}

\end{document}
